% Part: first-order-logic
% Chapter: introduction
% Section: models-theories

\documentclass[../../../include/open-logic-section]{subfiles}

\begin{document}

\olfileid{fol}{int}{mod}

\olsection{نمونے تے نظریے}

جدوں پہلے درجے دی منطق دی نحو تے معنویات
تعریف ہو جان، اسیں !!{structure}s تے معنوی
تصورات دیاں خاصیتاں دی جانچ شروع کر سکدے
آں۔ اسیں !!{derivation} نظام وی تعریف
کر سکدے آں تے انہاں دی پڑتال وی کر
سکدے آں۔ !!{sentence}s دے کسے سیٹ
لئی اسیں پچھ سکدے آں: کیہڑیاں
!!{structure}s ایس سیٹ دے سارے
!!{sentence}s نوں سچا کردیاں نیں؟
جے !!{sentence}s دا سیٹ~$\Gamma$
دتا ہووے، تاں اوہناں نوں پورا کرن والی
!!a{structure}~$\Struct{M}$ نوں
\emph{$\Gamma$ دا نمونہ} آکھدے آں۔
اسیں~$\Gamma$ توں شروع کر کے
اوہدے نمونے لبھن دی کوشش کر سکدے
آں---اوہ کیہو جیہے ہوندے نیں؟
اوہناں دا سائز کِنّا وڈا یا چھوٹا
ہونا لازمی اے؟ پر اسیں کسے اک
!!{structure} یا !!{structure}s دے
اک مجموعے توں شروع کر کے ایہ وی
پچھ سکدے آں: انہاں وچ کیہڑے
!!{sentence}s سچے نیں؟ کی کوئی
!!{sentence}s ایہناں !!{structure}s
دی \emph{عین شناخت} ایس معنی وچ
کردے نیں کہ اوہ تے صرف اوہی
اوہناں وچ سچے ہوندے نیں؟
ایہ سوال \emph{نمونیاں دے نظریے}
دا موضوع نیں۔ ایہی سوال
\emph{مسلّمی طریقے} دی بنیاد
وی نیں: !!{structure}s دے اک
مجموعے نوں کسے نظریے دے
مسلّماں، یعنی !!{sentence}s
دے سیٹ، نال بیان کرنا۔
ایہ ایس مشاہدے نال ممکن اے
کہ مسلّماں توں پہلے درجے دی
منطق وچ منطقی طور تے نکلن
والے عین اوہی !!{sentence}s
مسلّماں دے سارے نمونیاں وچ
سچے ہوندے نیں۔

اک بہت سادہ مثال لئی پیش ترتیب
ویکھو۔ پیش ترتیب اوہ تعلق~$R$
اے جیہڑا کسے سیٹ~$A$ اُتے
انعکاسی وی اے تے متعدی وی۔
% PNB-INT-006: غیر خالی دائرہ تے معیاری زبان دی ضمنی روایت واضح؛نواں source-error دعویٰ نہیں۔
اک غیر خالی سیٹ~$A$، جیہدے اُتے دو
تھاواں والا تعلق
$R \subseteq A \times A$
ہووے، اوہی کچھ اے جیہڑا
سانوں اک دو تھاواں والی
تعلق-علامت~$\Obj P$ والی
پہلے درجے دی زبان لئی
!!a{structure} دین لئی چاہیدا
اے: اسیں $\Domain{M} = A$
تے $\Assign{\Obj P}{M} = R$
رکھاں گے۔ چونکہ $R$ اک
پیش ترتیب اے، ایہ انعکاسی
تے متعدی اے؛ اسیں پہلے
درجے دی منطق دے
!!{sentence}s دا اک
سیٹ~$\Gamma$ لبھ سکدے
آں جیہڑا ایہ گل بیان کردا اے:
\begin{align*}
  & \lforall[\Obj v_0][\Atom{\Obj P}{\Obj v_0,\Obj v_0}]\\
  & \lforall[\Obj v_0][\lforall[\Obj v_1][\lforall[\Obj v_2][((\Atom{\Obj P}{\Obj v_0,\Obj v_1} \land \Atom{\Obj P}{\Obj v_1,\Obj v_2}) \lif \Atom{\Obj P}{\Obj v_0,\Obj v_2})]]]
\end{align*}
ایہ !!{sentence}s بس ”ہر~$x$ لئی
$Rxx$“ ($R$ انعکاسی اے)
تے ”جدوں $Rxy$ تے $Ryz$
ہون تاں $Rxz$ وی“ ($R$
متعدی اے) دیاں علامتی
لکھتاں نیں۔ اسیں ویکھدے
آں کہ !!a{structure}~$\Struct{M}$
ایہناں دو !!{sentence}s دے
سیٹ~$\Gamma$ دا نمونہ
اُدوں تے صرف اُدوں اے
جدوں $R$ (یعنی
$\Assign{\Obj P}{M}$)،
~$A$ اُتے پیش ترتیب
اے (یعنی $\Domain{M}$)۔
دوجے لفظاں وچ، $\Gamma$
دے نمونے عین پیش ترتیباں
نیں۔ ساریاں پیش ترتیباں
دی کوئی وی خاصیت جیہڑی
صرف~$\Obj P$ نوں
!!{predicate} منن والی
پہلے درجے دی زبان وچ
بیان ہو سکدی ہووے
(جیویں اُتّے انعکاسیت
تے تعدیت)، اوہ
~$\Gamma$ دے دو
!!{sentence}s توں منطقی
طور تے نکلدی اے،
تے اُلٹا وی۔ ایس لئی
اسیں~$\Gamma$ دے
نمونیاں بارے جو کجھ
ثابت کردے آں، اوہ
ساریاں پیش ترتیباں
بارے ثابت کردے آں۔

ہر خاص نظریے تے اوہدے
نمونیاں دے طبقے لئی
(جیویں $\Gamma$ تے
ساریاں پیش ترتیباں)،
ایہ دلچسپ سوال ہون
گے کہ مطابق پہلے درجے
دی زبان وچ کیہ بیان
ہو سکدا اے تے کیہ
نہیں۔ !!{structure}s
دیاں کجھ خاصیتاں
ایہوجیہیاں نیں جیہڑیاں
ہر زبان تے نمونیاں دے
ہر طبقے لئی اہم نیں،
خاص کر !!{domain}
دے سائز نال متعلق
خاصیتاں۔ ایتھے برابری والی،زیادہ توں زیادہ قابلِشمار
علامتی ذخیرے دی معیاری زبان تے اوہدیاں ذیلی زباناں مراد نیں۔ مثال لئی،
!!{domain} وچ عین
~$n$~!!{element}s
ہون دی گل کسے وی
~$n \in \PosInt$
لئی ہمیشہ بیان
کیتی جا سکدی اے۔
لامتناہی !!{sentence}s
دا اک سیٹ ورت کے
ایہ وی بیان کیتا
جا سکدا اے کہ
!!{domain} لامتناہی
اے۔ پر ایہ بیان
نہیں کیتا جا سکدا
کہ دائرہ محدود اے
یا دائرہ
!!{nonenumerable}
اے۔ پہلے درجے دیاں
زباناں دیاں حدّاں
بارے ایہ نتیجے
محدود گواہی تے
لووینہائم--سکولم
دے قضیاں توں نکلدے
نیں۔

\end{document}
